\chapter{Уравнения, неравенства и системы с параметром}\label{chap:exam-parameters}\markright{Глава 26}
\section{Параметр и особые значения}
\begin{definition}\term{Параметр} (parameter) — величина, которую при решении относительно неизвестного считают фиксированной, а затем исследуют зависимость ответа от её значения.\end{definition}
Задача «решить уравнение с параметром» требует описать корни при всех допустимых значениях параметра. Задача «найти значения параметра с данным числом решений» требует классифицировать эти значения, включая границы. Всюду считаем различные действительные корни; кратный корень считается одним, если специально не сказано иначе. Для системы считаем различные упорядоченные пары.
\begin{definition}\term{Особое значение параметра} (exceptional parameter value) в выбранном способе решения обращает в нуль коэффициент или делитель, меняет ОДЗ либо делает отдельный этап преобразования неприменимым. Такое значение рассматривают до деления.\end{definition}
\begin{theorem}[Линейное вырождение]
Уравнение \(A(a)x=B(a)\) при фиксированном \(a\) имеет один корень \(\frac{B(a)}{A(a)}\), если \(A(a)\ne0\); не имеет корней, если \(A(a)=0,B(a)\ne0\); верно для всех \(x\in\mathbb R\), если \(A(a)=B(a)=0\). Если исходная задача имела дополнительную ОДЗ, результат пересекают с ней.
\end{theorem}
\begin{proof}При ненулевом коэффициенте деление обратимо. При нулевом коэффициенте левая часть равна нулю независимо от \(x\), что даёт два оставшихся случая. Дополнительные ограничения остаются необходимыми по определению решения.\end{proof}
\par\addvspace{2pt}\begin{lrbox}{\CatFlowBox}\begin{minipage}[t]{\linewidth}
\setlength{\parskip}{2pt}\setlength{\columnsep}{4.5mm}\setlength{\intextsep}{0pt}
\begin{wrapfigure}[13]{r}{54mm}
\includegraphics[width=54mm,height=58mm,keepaspectratio]{images/unique/art-236.png}
\label{catflow:236}
\end{wrapfigure}
Параметр выбирают один раз для всего уравнения. После этого неизвестное может принимать разные значения, но коэффициенты уже зафиксированы. Поэтому сначала полезно выписать те значения параметра, при которых меняется степень уравнения, исчезает знаменатель или сливаются две границы. В каждом таком случае исходную запись проверяют отдельно. На остальных значениях можно применять общий способ и затем исследовать полученный ответ. Такая организация не требует заранее угадывать все корни. Она требует только понимать, какие действия обратимы и при каких условиях. Итоговая классификация должна покрывать всю прямую параметра без пропусков и пересечений несовместимых случаев. Для границ нужно своё объяснение, даже если они выглядят продолжением соседнего промежутка.
\par\WFclear\end{minipage}\end{lrbox}%
\ifdim\dimexpr\ht\CatFlowBox+\dp\CatFlowBox\relax<61mm\dp\CatFlowBox=\dimexpr61mm-\ht\CatFlowBox\relax\fi%
\noindent\usebox{\CatFlowBox}\par\addvspace{2pt}
\begin{example}Уравнение \((a^2-1)x=a-1\) при \(a=1\) превращается в \(0=0\), поэтому подходит любое \(x\). При \(a=-1\) получаем \(0=-2\), решений нет. Для \(a\ne\pm1\) деление даёт \(x=\frac1{a+1}\). Сокращение на \(a-1\) до разбора \(a=1\) потеряло бы целый случай.\end{example}

\section{Квадратное уравнение: дискриминант и Виет}
В уравнении \(A(a)x^2+B(a)x+C(a)=0\) сначала проверяют \(A(a)=0\). Для \(A(a)\ne0\) вычисляют \(D=B^2-4AC\): при \(D<0\) корней нет, при \(D=0\) один, при \(D>0\) два различных. Условие \(D=0\) не описывает линейные случаи, поэтому само по себе не отвечает на вопрос о единственном корне.
\begin{example}Уравнение \(ax^2-2x+1=0\) при \(a=0\) линейно и имеет один корень \(x=\frac12\). При \(a\ne0\) дискриминант равен \(4(1-a)\). Единственный корень получается также при \(a=1\), когда \((x-1)^2=0\). Искомые параметры для одного корня: \(a=0,1\).\end{example}
\begin{theorem}[Знаки двух действительных корней]
Пусть \(q(x)=x^2-sx+p\) имеет действительные корни с учётом кратности, то есть \(D=s^2-4p\geq0\). Оба корня положительны тогда и только тогда, когда \(s>0,p>0\); оба отрицательны тогда и только тогда, когда \(s<0,p>0\). Корни имеют противоположные знаки тогда и только тогда, когда \(p<0\). Для двух различных корней дополнительно нужно \(D>0\).
\end{theorem}
\begin{proof}По Виету сумма равна \(s\), произведение \(p\). Положительное произведение означает одинаковые ненулевые знаки, которые определяются суммой. Отрицательное произведение означает противоположные знаки и уже исключает совпадение корней. Обратные утверждения следуют из тех же суммы и произведения.\end{proof}
\begin{example}Найдём \(a\), при которых \(x^2-2ax+a=0\) имеет два различных положительных корня. Требуются одновременно \(4a(a-1)>0\), \(2a>0\), \(a>0\). Пересечение даёт \(a>1\). Условие одного дискриминанта допустило бы также \(a<0\), где произведение отрицательно и корни разных знаков.\end{example}

\section{Корни на отрезке}
\begin{theorem}[Оба корня на заданном отрезке]
Пусть \(l<u\), \(q(x)=x^2-sx+p\), \(v=\frac s2\), \(D=s^2-4p\). Оба корня с учётом кратности принадлежат \([l;u]\) тогда и только тогда, когда
\[
D\geq0,\qquad l\leq v\leq u,\qquad q(l)\geq0,\quad q(u)\geq0.
\]
Для двух различных корней заменяют \(D\geq0\) на \(D>0\). Чтобы оба корня лежали внутри \((l;u)\), берут \(l<v<u\), \(q(l)>0,q(u)>0\) и соответствующее условие на \(D\).
\end{theorem}
\begin{proof}Корни равны \(v-h,v+h\), где \(h=\frac{\sqrt D}2\geq0\), а \(q(x)=(x-v)^2-h^2\). Если оба корня на отрезке, то вершина между ними, \(v-l\geq h\), \(u-v\geq h\); отсюда условия на значения \(q\) в концах. Обратно, при \(l\leq v\leq u\) из \(q(l)\geq0\), \(q(u)\geq0\) следует \(v-l\geq h\), \(u-v\geq h\), так что оба корня на отрезке. Строгие условия дают строго внутреннее положение. Различность равносильна \(h>0\).\end{proof}
Для \(Ax^2+Bx+C\), \(A\ne0\), сначала делят на \(A\), чтобы получить приведённый многочлен. Знаки значений исходного многочлена при \(A<0\) противоположны знакам приведённого; применять критерий без нормировки нельзя.
\par\addvspace{2pt}\begin{lrbox}{\CatFlowBox}\begin{minipage}[t]{\linewidth}
\setlength{\parskip}{2pt}\setlength{\columnsep}{4.5mm}\setlength{\intextsep}{0pt}
\begin{wrapfigure}[11]{l}{54mm}
\includegraphics[width=54mm,height=58mm,keepaspectratio]{images/unique/art-237.png}
\label{catflow:237}
\end{wrapfigure}
Граница отрезка является частью условия задачи. Корень, совпавший с концом замкнутого отрезка, нужно включить, а с концом открытого промежутка — исключить. Если спрашивают ровно один корень на отрезке, недостаточно проверить смену знака в концах: двойной корень внутри может не менять знак, а два корня способны оказаться на самих концах. В таких задачах удобно найти корни явно, если формула проста, и для каждой серии выписать допустимый промежуток параметра. Затем выбирают значения, для которых подходит ровно одна серия, проверяя совпадения корней отдельно. Общий критерий для двух корней полезен как проверка, но не подменяет подсчёт различных решений в вопросе о единственности.
\par\WFclear\end{minipage}\end{lrbox}%
\ifdim\dimexpr\ht\CatFlowBox+\dp\CatFlowBox\relax<58mm\dp\CatFlowBox=\dimexpr58mm-\ht\CatFlowBox\relax\fi%
\noindent\usebox{\CatFlowBox}\par\addvspace{2pt}
\begin{example}Уравнение \(x^2-2ax+a^2-1=0\) имеет корни \(a-1,a+1\). Первый лежит на \([0;2]\) при \(1\leq a\leq3\), второй — при \(-1\leq a\leq1\). При \(a=1\) на отрезке находятся два различных корня \(0,2\). Поэтому ровно один корень на отрезке получается при \(a\in[-1;1)\cup(1;3]\).\end{example}

\section{Графические пересечения и касания}
\begin{definition}\term{Касание} (tangency) прямой и параболы в рассматриваемых задачах соответствует их единственной общей точке: после подстановки уравнения прямой в уравнение параболы получается квадратное уравнение с двойным корнем.\end{definition}
\begin{theorem}[Прямая и парабола]
Графики \(y=x^2\) и \(y=mx+b\) имеют две различные общие точки при \(m^2+4b>0\), одну при \(m^2+4b=0\) и ни одной при \(m^2+4b<0\). В случае единственной точки \(x=\frac m2\) разность высот равна \((x-\frac m2)^2\), поэтому прямая всюду лежит не выше параболы и касается её в этой точке.
\end{theorem}
\par\addvspace{2pt}\begin{lrbox}{\CatFlowBox}\begin{minipage}[t]{\linewidth}
\setlength{\parskip}{2pt}\setlength{\columnsep}{4.5mm}\setlength{\intextsep}{0pt}
\begin{wrapfigure}[9]{l}{54mm}
\includegraphics[width=54mm,height=58mm,keepaspectratio]{images/unique/art-239.png}
\label{catflow:239}
\end{wrapfigure}
\textit{Доказательство.} Уравнение абсцисс пересечений \(x^2-mx-b=0\) имеет дискриминант \(m^2+4b\). При нулевом дискриминанте \(b=-\frac{m^2}4\), и выделение квадрата даёт указанную разность.\hfill\(\blacksquare\)
\par\WFclear\end{minipage}\end{lrbox}%
\ifdim\dimexpr\ht\CatFlowBox+\dp\CatFlowBox\relax<58mm\dp\CatFlowBox=\dimexpr58mm-\ht\CatFlowBox\relax\fi%
\noindent\usebox{\CatFlowBox}\par\addvspace{2pt}

\begin{example}Для \(y=x^2\), \(y=2x+a\) дискриминант равен \(4(1+a)\). При \(a<-1\) пересечений нет, при \(a=-1\) одна точка \((1,1)\), при \(a>-1\) две. Так значение касания разделяет два разных количества решений.\end{example}
\par\addvspace{2pt}\begin{lrbox}{\CatFlowBox}\begin{minipage}[t]{\linewidth}
\setlength{\parskip}{2pt}\setlength{\columnsep}{4.5mm}\setlength{\intextsep}{0pt}
\begin{wrapfigure}[11]{r}{54mm}
\includegraphics[width=54mm,height=58mm,keepaspectratio]{images/unique/art-238.png}
\label{catflow:238}
\end{wrapfigure}
Графическое решение должно опираться на точные условия. Рисунок помогает обнаружить касание, пересечение границы или появление новой ветви, но приблизительная картинка не доказывает значение параметра. Для прямой и параболы условие касания проверяется дискриминантом или выделением полного квадрата. Для модуля вершины и линейные участки позволяют записать точные формулы без производной. Для системы каждое пересечение должно давать допустимую пару координат. После нахождения предполагаемых критических значений проверяют сами границы и промежутки между ними, используя полученные формулы или доказанную монотонность. Тогда график становится средством организации рассуждения: он подсказывает случаи, а алгебра объясняет, почему список случаев полный и сколько решений возникает в каждом из них.
\par\WFclear\end{minipage}\end{lrbox}%
\ifdim\dimexpr\ht\CatFlowBox+\dp\CatFlowBox\relax<58mm\dp\CatFlowBox=\dimexpr58mm-\ht\CatFlowBox\relax\fi%
\noindent\usebox{\CatFlowBox}\par\addvspace{2pt}
\section{Модуль, корень и число решений}
Для \(|x-c|=a\) при \(a<0\) корней нет, при \(a=0\) один, при \(a>0\) два. Несколько модулей раскрывают на промежутках с границами в их нулях; при зависимых от параметра границах сначала устанавливают их порядок.
\begin{example}Для \(|x-1|+|x+1|=a\) сумма равна \(2\) на \([-1;1]\), \(-2x\) слева и \(2x\) справа. Поэтому при \(a<2\) решений нет; при \(a=2\) их бесконечно много — весь отрезок; при \(a>2\) ровно два: \(x=\pm \frac a2\). Две совпавшие формулы корня не объяснили бы случай целого отрезка, поэтому плоский участок проверяем отдельно.\end{example}
\begin{example}Уравнение \(x^4-2ax^2+1=0\) после \(t=x^2\geq0\) даёт \(t^2-2at+1=0\). При \(a>1\) два различных положительных корня \(t=a\pm\sqrt{a^2-1}\) дают четыре корня \(x\). При \(a=1\) одно положительное значение \(t=1\) даёт два корня. При \(a<1\) положительных корней \(t\) нет: либо дискриминант отрицателен, либо при \(a\leq-1\) оба значения неположительны и фактически отрицательны, так как произведение равно \(1\). Значит, действительных корней \(x\) нет.\end{example}
При \(\sqrt{F(x)}=G(x,a)\) сначала проверяют \(F(x)\geq0\), \(G(x,a)\geq0\); только затем возводят в квадрат. В логарифмической задаче ограничения на аргументы и основания тоже зависят от параметра. Кандидаты, совпавшие с запрещённой точкой, не считаются корнями.

\section{Неравенства с параметром}
При умножении или делении неравенства на выражение от параметра его знак определяет направление сравнения. Нулевой случай проверяют отдельно. В методе интервалов положение нулей и запрещённых точек может менять порядок при изменении параметра; прежде чем строить таблицу, находят параметры совпадения границ.
\begin{example}Решим \((x-a)(x-2)\leq0\) при всех \(a\). Если \(a<2\), произведение неположительно на \([a;2]\). Если \(a>2\), ответ \([2;a]\). При \(a=2\) получаем \((x-2)^2\leq0\), поэтому подходит только \(x=2\). Во всех случаях можно написать \([\min(a,2);\max(a,2)]\), но такая запись должна включать объяснение случая совпавших корней.\end{example}

\section{Системы с параметром}
\begin{theorem}[Пересечение прямой с окружностью]
Для \(R>0\) система \(x+y=a\), \(x^2+y^2=R^2\) имеет две различные пары при \(a^2<2R^2\), одну при равенстве и ни одной при \(a^2>2R^2\).
\end{theorem}
\begin{proof}Из \((x+y)^2+(x-y)^2=2(x^2+y^2)\) получаем \((x-y)^2=2R^2-a^2\). Если правая часть положительна, два её квадратных корня дают две различные пары через \(x=\frac{a+(x-y)}2\), \(y=\frac{a-(x-y)}2\). При нуле получается одна пара, при отрицательном числе ни одной. Все преобразования обратимы.\end{proof}
В линейной системе сначала ищут параметры, при которых обращается в нуль определитель. Для этих значений исходные уравнения могут оказаться совместными и зависимыми или противоречивыми. В нелинейной системе подстановка обычно даёт одно уравнение, но количество его корней совпадает с количеством пар только тогда, когда каждому оставшемуся корню соответствует ровно одно допустимое значение второй неизвестной.

\relax

\par\addvspace{2pt}\begin{lrbox}{\CatFlowBox}\begin{minipage}[t]{\linewidth}
\setlength{\parskip}{2pt}\setlength{\columnsep}{4.5mm}\setlength{\intextsep}{0pt}
\begin{wrapfigure}[13]{r}{54mm}
\includegraphics[width=54mm,height=58mm,keepaspectratio]{images/unique/art-264.png}
\label{catflow:264}
\end{wrapfigure}
При подсчёте решений системы важно считать различные пары координат, а не способы получить их из уравнений. После разложения на множители одна и та же точка может лежать сразу на двух ветвях. Тогда простое сложение количеств пересечений завышает ответ. Отдельно проверяют значения параметра, при которых ветви совпадают, касаются или теряют часть области определения. Критическое значение не обязательно входит в ответ: в нём нужно заново подставить параметр в исходную систему. Чертёж показывает возможные случаи, а равенства и неравенства устанавливают их точные границы. Следующие примеры соединяют геометрический подсчёт, ограничения и проверку общих точек.
\par\WFclear\end{minipage}\end{lrbox}%
\ifdim\dimexpr\ht\CatFlowBox+\dp\CatFlowBox\relax<61mm\dp\CatFlowBox=\dimexpr61mm-\ht\CatFlowBox\relax\fi%
\noindent\usebox{\CatFlowBox}\par\addvspace{2pt}
\relax
\section{Несколько ветвей и подсчёт различных решений}
Равенство произведения нулю превращает задачу в совокупность. Решения разных ветвей могут совпасть: такая точка считается один раз. Полезно сначала выписать решения каждой ветви, затем найти их пересечения и только после этого считать различные точки.
\begin{example}[Прямая пересекает объединение двух кривых]\label{ex:review-ege-exam-parameters-1}
Найдём все \(a\), при которых система
\[
(y-x^2)(y-2x-1)=0,\qquad y=a
\]
имеет ровно два решения. Первая ветвь \(y=x^2\) даёт при \(a<0\) ноль точек, при \(a=0\) одну, при \(a>0\) две. Вторая ветвь \(y=2x+1\) даёт одну точку при любом \(a\): \(x=\frac{a-1}{2}\).
При \(a<0\) всего одна точка. При \(a=0\) точки \((0,0)\), \((-\frac12,0)\) различны, поэтому их две. При \(a>0\) предварительно получаем три точки. Две записи задают одну точку тогда и только тогда, когда
\[
a=\left(\frac{a-1}{2}\right)^2,
\quad a^2-6a+1=0,
\quad a=3\pm2\sqrt2.
\]
Оба значения положительны. В каждом случае одна точка прямой уже лежит на параболе, поэтому остаются две различные пары. Ответ: \(a\in\{0,3-2\sqrt2,3+2\sqrt2\}\). Сложение чисел решений без проверки совпадений потеряло бы два значения параметра.
\end{example}

\begin{theorem}[Единственная общая точка двух окружностей]
Пусть радиусы двух окружностей \(R,r\) положительны, а расстояние между центрами равно \(d\). При \(d>0\) ровно одна общая точка существует тогда и только тогда, когда \(d=R+r\) или \(d=|R-r|\). При \(d=0\) общих точек нет, если \(R\ne r\), и бесконечно много, если \(R=r\).
\end{theorem}
\begin{proof}
При \(d>0\) направим ось \(x\) вдоль центров: окружности имеют уравнения \(x^2+y^2=R^2\), \((x-d)^2+y^2=r^2\). Вычитание даёт единственную возможную абсциссу \(x_0=\frac{d^2+R^2-r^2}{2d}\). Если \(R^2-x_0^2>0\), возможны два противоположных ненулевых значения \(y\); если равно нулю — одно; если отрицательно — ни одного. Условие \(x_0=\pm R\) равносильно \((d-R)^2=r^2\) или \((d+R)^2=r^2\), то есть указанным двум условиям на положительные радиусы. При общем центре сравниваем радиусы непосредственно.
\end{proof}
\begin{example}[Два вида касания]\label{ex:review-ege-exam-parameters-2}
В системе \(x^2+y^2=4\), \((x-5)^2+y^2=a^2\) требуется \(a>0\) и единственное решение. Расстояние центров \(5\), радиусы \(2,a\). Внешнее касание даёт \(5=2+a\), то есть \(a=3\); внутреннее — \(5=|2-a|\), откуда с учётом положительности \(a=7\). Ответ: \(3,7\). Условие о положительном радиусе исключает вырожденную окружность-точку.
\end{example}

\section{Параметр и корень на ограниченном промежутке}
\begin{example}[Монотонность заменяет длинное возведение в квадрат]\label{ex:review-ege-exam-parameters-3}
Рассмотрим \(\sqrt{x}=a-x\) только при \(0\le x\le9\). Равносильная запись \(a=x+\sqrt{x}\) показывает: правая часть строго возрастает от \(0\) до \(12\). Поэтому ровно один корень на данном отрезке есть при \(0\le a\le12\), вне этого диапазона корней нет. Если нужен сам корень, положим \(t=\sqrt{x}\ge0\). Тогда \(t^2+t=a\),
\[
t=\frac{\sqrt{1+4a}-1}{2},\qquad
x=\left(\frac{\sqrt{1+4a}-1}{2}\right)^2.
\]
Отрицательный корень уравнения для \(t\) не подходит. Проверка концов \(a=0,12\) даёт \(x=0,9\), которые включены в условие.
\end{example}
Для нелинейной системы число корней после подстановки совпадает с числом пар лишь при однозначном восстановлении второй координаты. Если используются симметрия или квадратный корень \(y=\pm\sqrt{1-x^2}\), обе ветви проверяют с их концами. Совпадение двух формул в одной точке не создаёт второго решения.

\section{Многочлен с параметром: вырождение и совпадение корней}
\begin{example}[Сначала факторизация, затем число корней]\label{ex:review-ege-exam-parameters-4}
При каких \(a\) уравнение
\[
x\bigl(ax^3-3x^2+a(a^2-1)x-3(a^2-1)\bigr)
=2\bigl(ax^3-3x^2+a(a^2-1)x-3(a^2-1)\bigr)
\]
имеет ровно два различных действительных корня? Переносим всё в одну часть и группируем:
\[
(x-2)(ax-3)(x^2+a^2-1)=0.
\]
Корень \(x=2\) есть всегда. При \(a=0\) второй множитель постоянный ненулевой; третий даёт \(\pm1\), итого три корня. При \(0<|a|<1\) третий множитель даёт два различных корня внутри \((-1;1)\), а \(\frac3a\) по модулю больше \(3\), так что вместе с \(2\) получаем четыре разных корня. При \(a=\pm1\) третий множитель даёт \(0\); корни \(2,0,\frac3a\) различны, их три. При \(|a|>1\) третий множитель корней не имеет; остаются \(2\) и \(\frac3a\). Они совпадают только при \(a=\frac32\). Поэтому ответ
\[
a\in(-\infty;-1)\cup(1;\tfrac32)\cup(\tfrac32;+\infty).
\]
Проверка \(a=0\), где линейный множитель вырождается, и проверка совпадения корней необходимы: число множителей не равно числу различных решений.
\end{example}

\par\addvspace{2pt}\begin{lrbox}{\CatFlowBox}\begin{minipage}[t]{\linewidth}
\setlength{\parskip}{2pt}\setlength{\columnsep}{4.5mm}\setlength{\intextsep}{0pt}
\begin{wrapfigure}[12]{l}{54mm}
\includegraphics[width=54mm,height=58mm,keepaspectratio]{images/unique/art-280.png}
\label{catflow:280}
\end{wrapfigure}
Симметрия уравнения иногда определяет ответ раньше вычислений. Если замена знака одной координаты сохраняет всю систему, любое решение с ненулевой этой координатой имеет отличного от него партнёра. Поэтому единственную пару следует искать на оси симметрии. Полученные на оси значения параметра остаются лишь кандидатами: каждый нужно проверить на существование других точек вне оси. Для исключения лишних решений удобно собрать исходное равенство в сумму неотрицательных выражений. Тогда равенство нулю одновременно заставляет каждое слагаемое обращаться в нуль. Границы координат берут из второго уравнения системы, а знак каждого выражения проверяют на всём допустимом промежутке. Так устанавливают и необходимость, и достаточность ответа.
\par\WFclear\end{minipage}\end{lrbox}%
\ifdim\dimexpr\ht\CatFlowBox+\dp\CatFlowBox\relax<58mm\dp\CatFlowBox=\dimexpr58mm-\ht\CatFlowBox\relax\fi%
\noindent\usebox{\CatFlowBox}\par\addvspace{2pt}
\relax
\section{Симметрия системы и единственность}
Если оба уравнения системы сохраняются при замене \(x\) на \(-x\), то вместе с решением \((x,y)\) существует \((-x,y)\). Поэтому единственное решение обязательно имеет \(x=0\). Это необходимое условие: после нахождения кандидата параметра надо проверить, что других решений действительно нет.
\begin{example}[Экспонента, модуль и окружность]\label{ex:review-ege-last-full-system}
Разберём систему из открытого профильного варианта ФИПИ 2026 года, задания 18:
\[
\begin{cases}
2^{|x|+3}+7|x|+1=8y+7x^2+a,\\
x^2+y^2=1.
\end{cases}
\]
Оба уравнения чётны по \(x\). Если решение единственно, то \(x=0\), а окружность даёт \(y=1\) или \(y=-1\). Первое уравнение при \(x=0\) превращается в \(9=8y+a\). Кандидаты параметра: \(a=1\) и \(a=17\).
При \(a=17\) есть не только \((0,-1)\), но также \((1,0)\), \((-1,0)\): в последнем случае обе части первого равенства равны \(24\). Этот параметр не подходит.
При \(a=1\) положим \(u=|x|\). На окружности \(0\le u\le1\), \(y\le1\). Перенос членов даёт равносильное равенство
\[
8(2^u-1)+7u(1-u)+8(1-y)=0.
\]
Каждое слагаемое неотрицательно: \(2^u\ge1\), \(u(1-u)\ge0\), \(1-y\ge0\). Сумма равна нулю только тогда, когда все слагаемые нулевые. Из первого следует \(u=0\), из третьего \(y=1\). Эта пара удовлетворяет системе. Поэтому полный ответ \(\boxed{a=1}\), единственное решение \((0,1)\).
Здесь экспоненту не пришлось обращать: симметрия ограничила кандидатов, а знаки слагаемых доказали достаточность. Важно точно читать границы модуля: \(7|x|+1\) и \(7|x+1|\) задают разные задачи.
\end{example}

\section{Монотонная композиция и касание окружности}
\begin{example}[Нелинейная запись с простым геометрическим содержанием]\label{ex:review-ege-last-monotone}
Найдём \(a\), при которых система
\[
x^2+y^2=2,\qquad 2^{x+y}+3(x+y)=a
\]
имеет единственное решение. Функция \(H(s)=2^s+3s\) строго возрастает: при увеличении \(s\) строго возрастают оба слагаемых. Поэтому каждому \(a\) соответствует не более одного значения \(s=x+y\).
Из \((x-y)^2\ge0\) и \(x^2+y^2=2\) следует \(-2\le x+y\le2\). Значение \(s=2\) возможно только при \(x=y=1\), а \(s=-2\) только при \(x=y=-1\). Для любого \(-2<s<2\) из
\[
(x-y)^2=4-s^2>0,\qquad
x=\frac{s+(x-y)}2,\quad y=\frac{s-(x-y)}2
\]
получаются две разные пары. Значит, единственность бывает только на двух крайних уровнях: \(a=H(-2)=-\frac{23}{4}\), \(a=H(2)=10\).
Наличие экспоненты само по себе не делает задачу трудной: здесь она входит в строго монотонную функцию одной линейной комбинации. Для другой системы такую комбинацию или симметрию нужно сначала обосновать.
\end{example}
\section{Вырождение может совпасть с общим корнем ветвей}
В примере~\ref{ex:review-ege-exam-parameters-4} проверялись потеря линейного множителя при \(a=0\) и совпадение двух линейных корней. В следующем упражнении при \(a=0\) дополнительно совпадут постоянный корень и один из корней квадратного множителя. Поэтому особые значения параметра следует подставлять в целое произведение, а не только отдельно считать корни его множителей.
\section{Корень в произведении оставляет только часть ветви}
\begin{example}[Прямая и допустимая часть параболы]\label{ex:review-ege-final-parameters-A}
Для каждого действительного \(a\) найдём число различных пар в системе
\[
\begin{cases}
(y-x^2)\sqrt{x+2-y}=0,\\
y=a.
\end{cases}
\]
Исходная область определения: \(y\le x+2\). На ней произведение равно нулю тогда и только тогда, когда \(y=x^2\) или \(y=x+2\). Но первая ветвь сохраняет ограничение \(x^2\le x+2\), то есть \(-1\le x\le2\). Вторая ветвь целиком допустима: подкоренное выражение на ней равно нулю. Итак, пересекаем горизонталь \(y=a\) с отрезком параболы и со всей прямой, а не с двумя полными графиками.
Прямая даёт ровно одну пару \((a-2,a)\) при любом \(a\). На параболе при \(a<0\) точек нет; при \(a=0\) есть \((0,0)\); при \(a>0\) проверяем \(x=\pm\sqrt a\) на отрезке \([-1;2]\). Отрицательный корень допустим при \(0<a\le1\), положительный --- при \(0<a\le4\).
Ветви пересекаются при \(x^2=x+2\): это точки \((-1,1)\), \((2,4)\). Поэтому уровни \(a=1,4\) требуют устранить двойной счёт. При \(a=0\) точки \((0,0),(-2,0)\) различны. При \(0<a<1\) имеем три точки. При \(a=1\) одна пара записана дважды, и остаются две. При \(1<a<4\) остаются положительный корень параболы и точка прямой; они различны. При \(a=4\) эти две записи совпадают, а выше парабола точек не даёт.
\[
\begin{array}{c|ccccc}
a&(-\infty;0)&\{0\}&(0;1)&[1;4)&[4;+\infty)\\ \hline
\text{число пар}&1&2&3&2&1
\end{array}
\]
Следовательно, ровно две пары получаются при \(a\in\{0\}\cup[1;4)\). Ограничения ветви нашли до подсчёта, а совпадения проверили после него. Если сначала заменить допустимую часть всей параболой, исчезнет правильный ответ для \(1<a<4\).
\end{example}

\section{Сначала корни замены, затем число их прообразов}
\par\addvspace{2pt}\begin{lrbox}{\CatFlowBox}\begin{minipage}[t]{\linewidth}
\setlength{\parskip}{2pt}\setlength{\columnsep}{4.5mm}\setlength{\intextsep}{0pt}
\begin{wrapfigure}[13]{l}{54mm}
\includegraphics[width=54mm,height=58mm,keepaspectratio]{images/unique/art-287.png}
\label{catflow:287}
\end{wrapfigure}
Замена переменной разбивает задачу на два разных подсчёта. Сначала находят допустимые значения новой переменной, затем для каждого значения считают исходные решения. Даже два различных корня квадратного уравнения по новой переменной могут не дать двух корней исходного: один уровень может быть недостижим, другой может соответствовать целому отрезку. У суммы двух модулей минимальный уровень часто достигается между их вершинами. При совпадении этих вершин отрезок вырождается в точку, и число решений меняется ещё раз. Поэтому отдельно отмечают совпадение корней замены, достижимость каждого уровня и его положение относительно минимального значения исходного выражения.
\par\WFclear\end{minipage}\end{lrbox}%
\ifdim\dimexpr\ht\CatFlowBox+\dp\CatFlowBox\relax<61mm\dp\CatFlowBox=\dimexpr61mm-\ht\CatFlowBox\relax\fi%
\noindent\usebox{\CatFlowBox}\par\addvspace{2pt}
\relax

\begin{example}[Квадратное уравнение от суммы модулей]\label{ex:review-ege-final-parameters-B}
Для каждого действительного \(a\) определим число различных корней уравнения
\[
\bigl(|x-a|+|x+a|\bigr)^2-4\bigl(|x-a|+|x+a|\bigr)+4-a^2=0.
\]
Все выражения определены при любых действительных \(x,a\). Положим \(r=|a|\ge0\) и \(t=|x-a|+|x+a|\). Раскрытие модулей слева, между точками \(-r,r\) и справа от них даёт
\[
t=2\max(|x|,r)=
\begin{cases}2r,&|x|\le r,\\2|x|,&|x|>r.\end{cases}
\]
Поэтому уровень \(t<2r\) не имеет прообразов; уровень \(t>2r\) даёт ровно два корня \(x=\pm\frac t2\). На уровне \(t=2r\) решения составляют весь отрезок \([-r;r]\). При \(r=0\) этот отрезок --- одна точка, а не бесконечное множество.
Уравнение по \(t\) имеет вид \((t-2)^2=r^2\), его уровни \(t_-=2-r\), \(t_+=2+r\). При \(r=0\) уровни совпадают: \(t=2\), откуда только два корня \(x=\pm1\). Нельзя дважды учитывать корень квадратного уравнения.
При \(r>0\) сравним оба уровня с \(2r\). Нижний выше плоского уровня ровно при \(r<\frac23\), совпадает с ним при \(r=\frac23\). Верхний выше него при \(r<2\), совпадает при \(r=2\). Получаем полный ответ:
\[
\begin{array}{c|ccccc}
r&\{0\}&(0;\frac23)&\{\frac23,2\}&(\frac23;2)&(2;+\infty)\\ \hline
\text{число корней }x&2&4&\text{бесконечно много}&2&0
\end{array}
\]
При \(r=\frac23\) множество корней равно \([-\frac23;\frac23]\cup\{-\frac43,\frac43\}\); при \(r=2\) оно равно \([-2;2]\). В остальных случаях допустимый уровень \(t_\pm>2r\) даёт \(\pm\frac{2\pm r}{2}\), причём внутренний знак выбирают по уровню. Проверка плоского участка и совпадения уровней при нулевом параметре необходима: количество значений \(t\) само по себе не равно количеству \(x\).
\end{example}

\section{Единственный корень и отсутствие корней --- разные случаи}
\par\addvspace{2pt}\begin{lrbox}{\CatFlowBox}\begin{minipage}[t]{\linewidth}
\setlength{\parskip}{2pt}\setlength{\columnsep}{4.5mm}\setlength{\intextsep}{0pt}
\begin{wrapfigure}[12]{r}{54mm}
\includegraphics[width=54mm,height=58mm,keepaspectratio]{images/unique/art-288.png}
\label{catflow:288}
\end{wrapfigure}
Если замена знака неизвестной сохраняет уравнение и его область определения, каждый ненулевой корень имеет отличного от него партнёра. Значит, единственным корнем такого уравнения может быть только ноль. Однако подстановка нуля даёт лишь необходимые значения параметра: при них могут остаться и другие пары корней. Обратная ошибка возникает при больших значениях параметра, когда решений вовсе нет. Поэтому классификация должна различать ноль решений, единственный нулевой корень и симметричные пары ненулевых корней. Переход к модулю неизвестной уменьшает число случаев, но при возвращении к исходной переменной каждый положительный корень модуля даёт две разные точки.
\par\WFclear\end{minipage}\end{lrbox}%
\ifdim\dimexpr\ht\CatFlowBox+\dp\CatFlowBox\relax<59mm\dp\CatFlowBox=\dimexpr59mm-\ht\CatFlowBox\relax\fi%
\noindent\usebox{\CatFlowBox}\par\addvspace{2pt}
\relax

\begin{example}[Четвёртая степень и два симметричных модуля]\label{ex:review-ege-final-parameters-C}
Пусть \(d\ge0\). Найдём число различных действительных корней
\[
x^4+d^2=|x-d|+|x+d|.
\]
Уравнение чётно по \(x\). Положим \(u=|x|\ge0\). Как в предыдущем разборе, правая часть равна \(2\max(u,d)\), поэтому нормальная форма задачи
\[
u^4+d^2=2\max(u,d).
\]
Каждое положительное \(u\) даёт два корня \(x=\pm u\), а \(u=0\) --- один. Следовательно, единственный корень возможен только при \(u=0\); подстановка даёт \(d^2=2d\), то есть кандидаты \(d=0,2\). Однако для полного ответа надо проверить кандидатов и отличить отсутствие корней от наличия пары ненулевых корней.
Если \(0\le u\le d\), равенство имеет вид \(u^4=d(2-d)\). При \(d>2\) решений на этой части нет. При \(d=0\) получаем \(u=0\). При \(0<d\le2\) кандидат \(u=\sqrt[4]{d(2-d)}\) принадлежит этой части тогда и только тогда, когда
\[
d(2-d)\le d^4
\quad\Longleftrightarrow\quad
(d-1)(d^2+d+2)\ge0
\quad\Longleftrightarrow\quad d\ge1.
\]
Деление на \(d\) здесь допустимо, поскольку случай \(d=0\) уже выделен. Итак, при \(1\le d\le2\) эта часть даёт одно \(u\), положительное при \(d<2\) и нулевое при \(d=2\).
На части \(u\ge d\) рассмотрим многочлен \(F_d(u)=u^4-2u+d^2\). При \(d=0\) его неотрицательные корни \(0,\sqrt[3]2\). При \(0<d<1\) на \([d;1]\)
\[
F_d(u)\le u^4+u^2-2u=u(u-1)(u^2+u+2)\le0,
\]
причём равенство нулю невозможно: при \(u<1\) произведение строго отрицательно, а \(F_d(1)=d^2-1<0\). Для \(v>u\ge1\)
\[
F_d(v)-F_d(u)=(v-u)(v^3+v^2u+vu^2+u^3-2)>0.
\]
Значит, на \([1;+\infty)\) многочлен строго возрастает. Он отрицателен при \(1\) и положителен при \(2\), поэтому имеет ровно один корень между ними. Здесь существование следует из свойства промежуточного значения непрерывного многочлена; его строгое доказательство дано в главе \ref{chap:c28}. Эта ссылка нужна для существования корня, а не для оценки его числа; формулу четвёртой степени мы не используем.
При \(d=1\) на \(u\ge1\) единственный корень \(u=1\), уже полученный на первой части. При \(d>1\) имеем \(F_d(d)=d(d^3+d-2)>0\); далее многочлен строго возрастает, поэтому на \(u\ge d\) корней нет.
После объединения частей и удаления совпавших граничных записей:
\[
\begin{array}{c|cccc}
d&\{0\}&(0;2)&\{2\}&(2;+\infty)\\ \hline
\text{число корней }x&3&2&1&0
\end{array}
\]
При \(d=0\) это \(0,\pm\sqrt[3]2\). При \(d=2\) единственный корень \(0\). Поэтому условие «ровно один корень или нет корней» равносильно \(d\ge2\), но единственность отдельно равносильна \(d=2\). Например, если \(d=|2a-1|\), единственность равносильна \(a=-\frac12\) или \(a=\frac32\); для отсутствия корней нужны строгие неравенства \(a<-\frac12\) или \(a>\frac32\).
\end{example}

\Needspace{7\baselineskip}\section*{Задачи}
\begingroup\setlength{\parskip}{1pt}
\noindent\begin{minipage}[t]{\linewidth}\textbf{26.1.} Реши относительно \(x\) уравнение \((a-2)x=3\) при всех \(a\).\par\end{minipage}\par
\noindent\begin{minipage}[t]{\linewidth}\textbf{26.2.} Реши \((a-1)x=a-1\) при всех \(a\).\par\end{minipage}\par
\noindent\begin{minipage}[t]{\linewidth}\textbf{26.3.} Реши \((a^2-1)x=a-1\) при всех \(a\).\par\end{minipage}\par
\noindent\begin{minipage}[t]{\linewidth}\textbf{26.4.} При каких \(a\) уравнение \(ax^2-2x+1=0\) имеет ровно один корень?\par\end{minipage}\par
\noindent\begin{minipage}[t]{\linewidth}\textbf{26.5.} При каких \(a\) уравнение \((a-1)x^2-2(a-1)x+a=0\) имеет два различных корня?\par\end{minipage}\par
\par\addvspace{2pt}\begin{lrbox}{\CatFlowBox}\begin{minipage}[t]{\linewidth}
\setlength{\parskip}{2pt}\setlength{\columnsep}{4.5mm}\setlength{\intextsep}{0pt}
\begin{wrapfigure}[11]{r}{54mm}
\includegraphics[width=54mm,height=58mm,keepaspectratio]{images/unique/art-240.png}
\label{catflow:240}
\end{wrapfigure}
\textbf{26.6.}\starmark{} Найди \(a\), при которых \(x^2-2ax+a=0\) имеет два различных положительных корня.\par

\textbf{26.7.}\starmark{} При каких \(a\) оба различных корня \(x^2-(a+2)x+2a=0\) лежат в \((0;3)\)?\par

\textbf{26.8.} Когда корни \(x^2-2ax+a+2=0\) имеют противоположные знаки?\par

\textbf{26.9.} При каких \(a\) оба корня \(x^2-2ax+a^2-4=0\) лежат на \([0;5]\)?\par
\par\WFclear\end{minipage}\end{lrbox}%
\ifdim\dimexpr\ht\CatFlowBox+\dp\CatFlowBox\relax<58mm\dp\CatFlowBox=\dimexpr58mm-\ht\CatFlowBox\relax\fi%
\noindent\usebox{\CatFlowBox}\par\addvspace{2pt}

\noindent\begin{minipage}[t]{\linewidth}\textbf{26.10.}\starmark{} Когда оба корня с учётом кратности \(x^2-2ax+a=0\) лежат на \([0;1]\)?\par\end{minipage}\par
\noindent\begin{minipage}[t]{\linewidth}\textbf{26.11.}\starmark{} При каких \(a\) уравнение \(x^2-2ax+a=0\) имеет ровно один различный корень на \([0;1]\)?\par\end{minipage}\par
\noindent\begin{minipage}[t]{\linewidth}\textbf{26.12.}\starmark{} При каких \(a\) ровно один корень \(x^2-2ax+a^2-1=0\) принадлежит \([0;2]\)?\par\end{minipage}\par
\noindent\begin{minipage}[t]{\linewidth}\textbf{26.13.} Когда \(x^2-2ax+1>0\) для всех действительных \(x\)?\par\end{minipage}\par
\noindent\begin{minipage}[t]{\linewidth}\textbf{26.14.} Когда \(x^2-2ax+1\geq0\) для всех действительных \(x\)?\par\end{minipage}\par
\noindent\begin{minipage}[t]{\linewidth}\textbf{26.15.}\starmark{} Найди все \(a\), при которых \(ax^2+2x+1\geq0\) при любом \(x\).\par\end{minipage}\par
\noindent\begin{minipage}[t]{\linewidth}\textbf{26.16.} При каких \(a\) уравнение \(|x-2|=a\) имеет два корня? Найди их.\par\end{minipage}\par
\noindent\begin{minipage}[t]{\linewidth}\textbf{26.17.}\starmark{} Определи число решений \(|x-2|+|x+2|=a\) при всех \(a\).\par\end{minipage}\par
\noindent\begin{minipage}[t]{\linewidth}\textbf{26.18.}\starmark{} При каких \(a\) уравнение \(|x^2-1|=a\) имеет четыре различных корня?\par\end{minipage}\par
\noindent\begin{minipage}[t]{\linewidth}\textbf{26.19.}\starmark{} Определи число решений \(|x-a|=x+1\) при всех \(a\).\par\end{minipage}\par
\noindent\begin{minipage}[t]{\linewidth}\textbf{26.20.} Для какого \(a\) прямая \(y=2x+a\) касается параболы \(y=x^2\)? Найди общую точку.\par\end{minipage}\par
\noindent\begin{minipage}[t]{\linewidth}\textbf{26.21.}\starmark{} При каких \(a\) система \(x+y=a\), \(x^2+y^2=8\) имеет ровно два решения?\par\end{minipage}\par
\noindent\begin{minipage}[t]{\linewidth}\textbf{26.22.}\starmark{} Реши систему \(ax+y=1\), \(x+ay=1\) при всех \(a\).\par\end{minipage}\par
\noindent\begin{minipage}[t]{\linewidth}\textbf{26.23.} При каких \(a\) система \(y=x^2\), \(y=ax\) имеет ровно одно решение?\par\end{minipage}\par
\noindent\begin{minipage}[t]{\linewidth}\textbf{26.24.}\starmark{} Найди число действительных корней \(x^4-2ax^2+1=0\) при всех \(a\).\par\end{minipage}\par
\noindent\begin{minipage}[t]{\linewidth}\textbf{26.25.}\starmark{} При каких \(a\) уравнение \(\sqrt{x+1}=a-x\) имеет ровно один корень? Найди его.\par\end{minipage}\par
\noindent\begin{minipage}[t]{\linewidth}\textbf{26.26.}\starmark{} При каких \(a\) уравнение \(\log_2(x-a)+\log_2(x+a)=0\) имеет единственный корень?\par\end{minipage}\par
\noindent\begin{minipage}[t]{\linewidth}\textbf{26.27.} При каких \(a\) уравнение \(2^x=a-1\) имеет корень? Сколько их?\par\end{minipage}\par
\noindent\begin{minipage}[t]{\linewidth}\textbf{26.28.}\starmark{} При каких \(a\) система \(x+y=a\), \(xy=1\) имеет две различные упорядоченные пары решений?\par\end{minipage}\par
\noindent\begin{minipage}[t]{\linewidth}\textbf{26.29.} При каких \(a\) уравнение \(\dfrac{x-a}{x-1}=0\) имеет единственный корень?\par\end{minipage}\par
\noindent\begin{minipage}[t]{\linewidth}\textbf{26.30.}\starmark{} Определи число корней \(\dfrac a{x-1}=x-1\) при всех \(a\).\par\end{minipage}\par
\noindent\begin{minipage}[t]{\linewidth}\textbf{26.31.}\starmark{} При каких \(a\) уравнение \(|x|=ax+1\) имеет ровно один корень?\par\end{minipage}\par
\noindent\begin{minipage}[t]{\linewidth}\textbf{26.32.}\starmark{} Найди \(a\), при которых \(x^2-2(a+1)x+a^2=0\) имеет один корень меньше \(1\), а другой больше \(1\).\par\end{minipage}\par
\noindent\begin{minipage}[t]{\linewidth}\textbf{26.33.} Реши неравенство \((a-1)x\geq2\) при всех \(a\).\par\end{minipage}\par
\noindent\begin{minipage}[t]{\linewidth}\textbf{26.34.}\starmark{} Реши \(\dfrac{x-a}{x-1}\geq0\) при всех \(a\).\par\end{minipage}\par
\noindent\begin{minipage}[t]{\linewidth}\textbf{26.35.}\starmark{} Окружности заданы уравнениями \(x^2+y^2=4\), \((x-6)^2+y^2=a^2\), где \(a>0\). При каких \(a\) система имеет ровно одно решение?\par\end{minipage}\par
\noindent\begin{minipage}[t]{\linewidth}\textbf{26.36.}\starmark{} При каких \(a\) система \((y-x^2)(y-3x+2)=0\), \(y=a\) имеет ровно две различные пары решений?\par\end{minipage}\par
\noindent\begin{minipage}[t]{\linewidth}\textbf{26.37.}\starmark{} Найди все \(a\), при которых система \((y-x^2)(y-x)=0\), \(y=a\) имеет ровно два решения.\par\end{minipage}\par
\noindent\begin{minipage}[t]{\linewidth}\textbf{26.38.} Для всех \(a\) определи число корней \(\sqrt{x}=a-x\) на \([0;4]\). При существовании корня вырази его через \(a\).\par\end{minipage}\par
\noindent\begin{minipage}[t]{\linewidth}\textbf{26.39.}\starmark{} При каких \(a\) уравнение \(\sqrt{x^2-2x+1}=a-x\) имеет ровно один корень на \([0;3]\)?\par\end{minipage}\par
\noindent\begin{minipage}[t]{\linewidth}\textbf{26.40.}\starmark{} Для каждого \(a\) найди число различных решений системы \(x^2+y^2=1\), \(y=a-|x|\). В частности, когда решение единственно?\par\end{minipage}\par
\noindent\begin{minipage}[t]{\linewidth}\textbf{26.41.}\starmark{} При каких \(b\) система \(x^2+y^2=1\), \(3^{|x|+2}+5|x|+2=18y+5x^2+b\) имеет ровно одну пару решений? Найди эту пару.\par\end{minipage}\par
\noindent\begin{minipage}[t]{\linewidth}\textbf{26.42.}\starmark{} Найди все \(c\), при которых система \(x^2+y^2=1\), \(2^{|x|+2}+2|x|+3=4y+2x^2+c\) имеет единственное решение. Обоснуй достаточность найденного значения.\par\end{minipage}\par
\noindent\begin{minipage}[t]{\linewidth}\textbf{26.43.}\starmark{} Пусть
\[
Q_a(x)=ax^3-4x^2+a(a^2-1)x-4(a^2-1).
\]
Найди все действительные \(a\), при которых уравнение \(xQ_a(x)=-Q_a(x)\) имеет ровно два различных действительных корня. Сначала разложи \(Q_a\) на множители; отдельно исследуй \(a=0\), \(a=\pm1\) и совпадение линейных корней.\par\end{minipage}\par
\noindent\begin{minipage}[t]{\linewidth}\textbf{26.44.}\starmark{} Для всех действительных \(b\) определи число различных пар в системе \((y-x^2)\sqrt{2x+3-y}=0\), \(y=b\). При каких \(b\) пар ровно две? Сохрани ограничения обеих ветвей.\par\end{minipage}\par
\noindent\begin{minipage}[t]{\linewidth}\textbf{26.45.}\starmark{} Пусть \(S_b(x)=|x-b|+|x+b|\), где \(b\in\mathbb R\). Для всех \(b\) найди число различных действительных корней \(S_b(x)^2-6S_b(x)+9-b^2=0\). В случаях бесконечного множества выпиши его полностью.\par\end{minipage}\par
\noindent\begin{minipage}[t]{\linewidth}\textbf{26.46.}\starmark{} Для всех действительных \(b\) определи число различных корней \(x^4+(b+1)^2=|x-b-1|+|x+b+1|\). Отдельно укажи параметры единственности и параметры отсутствия корней.\par\end{minipage}\par
\noindent\begin{minipage}[t]{\linewidth}\textbf{26.47.}\starmark{} Пусть \(Q_a(x)=ax^3-2x^2+a(a^2-1)x-2(a^2-1)\). Для каждого действительного \(a\) определи число различных действительных корней \(xQ_a(x)=3Q_a(x)\). При каких \(a\) корней ровно два?\par\end{minipage}\par
\par\Needspace{12\baselineskip}\section*{Олимпиадные задачи}
\par\addvspace{2pt}\begin{lrbox}{\CatFlowBox}\begin{minipage}[t]{\linewidth}
\setlength{\parskip}{2pt}\setlength{\columnsep}{4.5mm}\setlength{\intextsep}{0pt}
\begin{wrapfigure}[11]{l}{54mm}
\includegraphics[width=54mm,height=58mm,keepaspectratio]{images/unique/art-309.png}
\label{catflow:309}
\end{wrapfigure}
\textbf{26.48.}\starmark{} {\small\itshape\href{https://mmo.mccme.ru/books/mmo-35-57.pdf}{ММО, 1945, I тур, 9–10 классы, № 3}; редакция условия.}\par При каких действительных \(a\) система \(x^2-y^2=0\), \((x-a)^2+y^2=1\) имеет ровно три или ровно два различных действительных решения?\par

\textbf{26.49.}\starmark{} {\small\itshape\href{https://www.imo-official.org/assets/documents/problems/1963/1963_eng.pdf}{IMO, 1963, № 1}; перевод.}\par Для каждого действительного \(p\) найди все действительные корни \[\sqrt{x^2-p}+2\sqrt{x^2-1}=x.\]\par
\par\WFclear\end{minipage}\end{lrbox}%
\ifdim\dimexpr\ht\CatFlowBox+\dp\CatFlowBox\relax<58mm\dp\CatFlowBox=\dimexpr58mm-\ht\CatFlowBox\relax\fi%
\noindent\usebox{\CatFlowBox}\par\addvspace{2pt}

\endgroup
